\chapter{Допуск обобщённого первопорядкового родителя
\texorpdfstring{$R_2$}{R2} на \texorpdfstring{$H_{15}$}{H15}}
\label{ch:v10-h15-r2-generalized-first-order-parent}

Предыдущий аудит выделил минимальную типизированную пару
\begin{equation}
 D_{R_2}:\qquad L_L-u_R,\quad Q_L-e_R,
\end{equation}
но строгий первый порядок запретил оба блока. Проверим, может ли обобщённая
внутренняя флуктуация породить эту пару из стандартного оператора
$D_{\rm adm}$, поддержанного на $Q_L-u_R$, $Q_L-d_R$, $L_L-e_R$.

Без условия первого порядка ковариантная флуктуация имеет вид
\begin{equation}
 D_A=D+A_{(1)}+\widetilde A_{(1)}+A_{(2)},\qquad
 A_{(2)}=\sum_j\widehat a_j[A_{(1)},\widehat b_j].
\end{equation}
Возьмём центральные знаковые унитарии, различающие две бимодульные
координаты,
\begin{equation}
 L=\operatorname{diag}(1,1,-1,-1,-1),\qquad
 R=\operatorname{diag}(1,-1,1,1,-1).
\end{equation}
Для $U=LR$ точное разложение $UDU^*=D+A_{(1)}+\widetilde
A_{(1)}+A_{(2)}$ даёт
\begin{equation}
 A_{(2)}[D_{\rm adm}]=0.
\end{equation}
Это не особенность выбранных знаков: умножение на представленные элементы
алгебры и взятие коммутаторов не создают матричный блок там, где исходный
$D$ имел нулевую опору.

Положительный контроль начинается с заранее вставленного $D_{R_2}$. Для
него двойной коммутатор имеет ранг $4$, а выбранная флуктуация даёт
\begin{equation}
 A_{(1)}[D_{R_2}]=-2D_{R_2},\qquad
 \widetilde A_{(1)}[D_{R_2}]=-2D_{R_2},\qquad
 A_{(2)}[D_{R_2}]=4D_{R_2}.
\end{equation}
Следовательно, при $D=D_{\rm adm}+\varepsilon D_{R_2}$ квадратичный ответ
равен $4\varepsilon D_{R_2}$: его опора и амплитуда наследуются от seed,
а не возникают самостоятельно.

Графовое следствие непосредственно затрагивает равномерный усилитель.
Без seed лапласиан сохраняет ранг/ядро $3/2$, то есть независимые кварковую
и лептонную амплитуды. После условной вставки двух $R_2$-рёбер получаем
$4/1$, но inherited seed имеет ранг $0$.

\begin{theorem}[Сохранение опоры обобщённой флуктуацией]
На фиксированном представлении $H_{15}$ линейные и квадратичная внутренние
флуктуации не создают отсутствующий кварк--лептонный блок. Для стандартного
первопорядкового seed имеем $A_{(2)}=0$; ненулевой $R_2$-ответ
пропорционален предварительно вставленному $D_{R_2}$.
\end{theorem}

\begin{nogocriterion}[Круговой $R_2$-родитель]
Калибровочная ковариантность формулы без первого порядка не является
происхождением нового Dirac-блока. Если $R_2$ присутствует только как seed
исходного $D$, обобщённая флуктуация одевает уже постулированный сектор и не
выводит ни его наличие, ни coupling, ни нормировку uniform amplifier.
\end{nogocriterion}

\begin{protocol}\begin{enumerate}
\item обобщённое ковариантное разложение: \textbf{точно выполнено};
\item $A_{(2)}$ на стандартной опоре: \textbf{$0$};
\item ядро графа без $R_2$: \textbf{$2$};
\item $R_2$-seed активирует $A_{(2)}$: \textbf{да, $4D_{R_2}$};
\item опора $A_{(2)}$ шире опоры seed: \textbf{нет};
\item условное ядро после seed: \textbf{$1$};
\item inherited $R_2$ seed: \textbf{$0/2$};
\item physical origin: \textbf{$0/3$};
\item следующий гейт: \textbf{аудит происхождения конечного $R_2$-seed}.
\end{enumerate}\end{protocol}

Машинный сертификат сохранён в
\path{s2t_v10_particle_wrinkle_dislocation_callias_equal_charge_twist_m4_h15_r2_generalized_first_order_parent_admission_gate_results.json}.

\begin{thebibliography}{9}\raggedright
\bibitem{v10r2-generalized}
A.~H.~Chamseddine, A.~Connes, W.~D.~van Suijlekom,
\emph{Inner Fluctuations in Noncommutative Geometry without the First Order Condition},
Journal of Geometry and Physics 73 (2013), 222--234; arXiv:1304.7583.
\end{thebibliography}